\vspace{0.8cm}

5. En una ciudad, el 40\%V de la población trabaja en el sector de servicios, 
el 35\% en el sector de la salud, y el resto en otros sectores. 
El 30\% de los trabajadores en el sector de servicios son mujeres, el 60\% 
en el sector sanidad y el 45 \% en otros sectores. ¿Cuál es la probabilidad 
de que un hombre elegido al azar trabaje en sanidad? 

\vspace{0.5cm}

Vamos a definir los eventos.

$S$: ``Trabaja en servicios''

$Sa$: ``Trabaja en sanidad''

$O$: ``Trabaja en otros sectores''

$M$: ``Es mujer''

$\bar{M}$: ``Es onvre''

\vspace{0.4cm}

Sabemos que la proporción de la población que trabaja en cada sector es
$$P(S)=0.40,\quad P(Sa)=0.35,\quad P(O)=0.25.$$

\vspace{0.2cm}

Además, las proporciones de mujeres dentro de cada sector son

$$P(M\mid S)=0.30\;\Rightarrow\;P(\bar{M}\mid S)=0.70$$

$$P(M\mid Sa)=0.60\;\Rightarrow\;P(\bar{M}\mid Sa)=0.40$$

$$P(M\mid O)=0.45\;\Rightarrow\;P(\bar{M}\mid O)=0.55$$

\vspace{0.2cm}

Queremos calcular la probabilidad de que un hombre elegido al azar trabaje en sanidad, es decir

$$P(Sa\mid \bar{M})=\frac{P(Sa\cap\bar{M})}{P(\bar{M})}$$

$$P(Sa\cap\bar{M})=P(Sa)\,P(\bar{M}\mid Sa)=0.35\times 0.40=0.14$$

\vspace{0.2cm}

Por el teorema de probabilidad total
$$P(\bar{M})=P(\bar{M}\mid S)P(S)+P(\bar{M}\mid Sa)P(Sa)+P(\bar{M}\mid O)P(O)$$

$$P(\bar{M})=(0.70)(0.40)+(0.40)(0.35)+(0.55)(0.25)=0.28+0.14+0.1375=0.5575$$

\vspace{0.2cm}

Por lo tanto,
$$P(Sa\mid \bar{M})=\frac{0.14}{0.5575}\approx 0.2511.$$

La probabilidad de que un hombre elegido al azar trabaje en sanidad es $\boxed{25.11\%}$

\QED

\vspace{0.8cm}

6. Un médico ha observado que el 40\% de sus pacientes fuma y de estos, 
el 75\% son hombres. Entre los que no fuman, el 60\% son mujeres. 
Calcula la probabilidad de que una mujer fume.

\vspace{0.5cm}

Vamos a definir los eventos. 

$F$: ''Fuma''

$\bar{F}$: ''No fuma''

$M$: ''Es mujer''

$H$: ''Es hombre''

\vspace{0.4cm}

Sabemos que la probabilidad de que el paciente fume es del 40\%, es decir
$P(F) = 0.40,\ P(\bar{F}) = 0.60$. 
\vspace{0.2cm}

Por otro lado, la probabilidad de que el paciente sea hombre dado que fume, 
es $P(H \mid F) = 0.75$, entonces $P(M \mid F) = 0.25$.
\vspace{0.2cm}

Además, tenemos que la probabilidad de que el paciente sea mujer, dado a que fume es$P(M \mid \bar{F}) = 0.60$, entonces $P(H \mid \bar{F}) = 0.40$
\vspace{0.2cm}

Queremos calcular entcontrar la probabilidad de que un paciente fume, dado a que es 
mujer, es decir:
\vspace{0.2cm}

$$P(F \mid M) = \frac{P(F \cap M)}{P(M)}$$

Entonces

$$P(F \cap M) = P(F)\, P(M \mid F) = 0.40 \times 0.25 = 0.10$$

\vspace{0.2cm}

Utilizando el teorema de probabilidad total:

$$P(M) = P(M \mid F)\, P(F) + P(M \mid \bar{F})\, P(\bar{F})$$

$$P(M) = (0.25)(0.40) + (0.60)(0.60) = 0.10 + 0.36 = 0.46$$

Por lo tanto
\vspace{0.2cm}

$$P(F \mid M) = \frac{0.10}{0.46} \approx 0.217$$

La probabilidad de que un paciente mujer fume es $\boxed{21.7\%}$

\QED

\vspace{0.8cm}

7. En una fábrica se embalan galletas, en 4 cadebas de montaje, $A_1$ 
(saludos a la página de nombre largo), $A_2, A_3 \text{ y } A_4$. 
El 35\% de la producción total se embala en la cadena $A_1$, 
y el 20\%, 24\% y 21\% en $A_2$, $A_3$ y $A_4$, respectivamente. 
Los datos indican que no se embalan correctamente un porcentaje 
pequeño de las cajas; el 1\% de $A_1$, el 3\% de $A_2$, el 2.5\% de $A_3$ 
y el 2\% de $A_4$. ¿Cuál es la probabilidad de que una caja elegida al azar 
de la producción total sea defectuosa?

\vspace{0.5cm}

Vamos a definir los eventos.

$A_1, A_2, A_3, A_4$: ``La caja se embala en la cadena $i$'', $i=1,2,3,4$.

$D$: ``La caja está defectuosa''

\vspace{0.4cm}

Sabemos que la proporción de producción en cada cadena es
$$P(A_1)=0.35,\quad P(A_2)=0.20,\quad P(A_3)=0.24,\quad P(A_4)=0.21$$

\vspace{0.2cm}

Además, la probabilidad de que una caja sea defectuosa en cada cadena es

$$P(D\mid A_1)=0.01,\quad P(D\mid A_2)=0.03,\quad P(D\mid A_3)=0.025,\quad P(D\mid A_4)=0.02$$

\vspace{0.2cm}

Aplicando el teorema de probabilidad total, la probabilidad de que una caja elegida al azar sea defectuosa es

$$P(D)=P(D\mid A_1)P(A_1)+P(D\mid A_2)P(A_2)+P(D\mid A_3)P(A_3)+P(D\mid A_4)P(A_4)$$

$$P(D) = (0.01)(0.35) + (0.03)(0.20) + (0.025)(0.24) + (0.02)(0.21)$$

$$P(D) = 0.0035 + 0.006 + 0.006 + 0.0042 = 0.0197.$$

\vspace{0.2cm}

Por lo tanto, la probabilidad de que una caja elegida al azar sea defectuosa es $\boxed{1.97\%}$.

\QED

\vspace{0.8cm}
